\section{Discussion}

\subsection{What the experiments establish}

The strongest physical result is an interaction between outage position, target statistic and reconstruction information. Missing-day counts are insufficient for H30, but peak overlap is not sufficient either. Explicit support can help transparent rules and restricted models; it is not the necessary component behind the best flexible score. The proposed representation does not establish its intended increment over strong target-specific diagnostics. This negative finding narrows, rather than invalidates, the value of a fit-for-use assessment.

For a fixed point estimate, a complete conditional distribution of H30 would already determine the probability of exceeding its tolerance. Window summaries can help estimate that distribution, but an additional decision signal cannot be assumed. Our direct uncertainty summaries and trajectories are approximations: their observed redundancy is empirical, not a proof that every support representation is useless. Window information may be more relevant to choosing a costly additional data source than to repeatedly summarizing one fixed estimate's risk.

The decomposition makes a separate point: a nonnegative term in a maximum-statistic identity need not describe a harmful process. Switching can compensate a damaged true peak. Centering hidden residuals is likewise diagnostic, not a field method. Interpretation must respect the final sign, intervention interaction and small improvement demonstrated from visible information.

\subsection{Why deployment conditions matter}

New-station tests distinguish ranking from calibration transport. A useful source score can miss its target elsewhere, even with intercontinental river-flow separation. Local calibration changes the problem and may help; it is not successful unchanged transfer. Future station history is legitimate retrospective archive information, not prospective information.

Stream-temperature work already shows that similar in-distribution predictive accuracy need not imply similar performance under environmental shift \cite{topp2023}. Our result concerns transport of an admission calibration, not a demonstration that one neural architecture better represents river processes. Both forms of generalization need tests matching the claimed deployment conditions.

Risk statements must match the independent unit and population. Overlap components prevent duplicate evidence but do not establish independent catchment climates. The formal rule is limited by that assumption and its group-expectation target. No certifiable threshold means insufficient evidence, not that every record is unusable. Nonzero empirical retention cannot be relabelled a guarantee. More masks cannot replace independent calibration groups.

\subsection{Magnitude, judgement and monitoring value}

An agency may need an H30 number, a hot-year label or an unbiased regional frequency. They incur different losses. A refused magnitude can establish hot conditions through a full observed window, but few labels were rescued here. Observations below threshold cannot establish cool conditions. Selected accuracy does not establish representativeness: unknown fractions, hot/cool refusal, coverage and original-pool frequency uncertainty remain outcomes.

The cost grid identifies where screening could help a stated decision. It does not estimate field-revisit costs, ecological losses or the harm of withholding a record. Those require prospective agency-specific work. Our tolerance and finite-history anomaly threshold are experimental choices, not biological standards. We recommend neither a universal archive cutoff nor an assumption that natural outages are harmless.

The main downstream no-screen reference uses the common priority reconstruction. Beating it does not establish a useful selective policy when declaring all answers unknown is cheaper. Any minimum information-return requirement belongs in a prospective utility specification, not a post-test constraint added to defeat that baseline. Scenario benefits also do not establish superiority over an optimally chosen physical reconstruction or information-acquisition policy.

\subsection{Limits and the next independent test}

Opened sources supported development and are not sealed confirmation. Even broader controls require complete training years and do not represent all short, intermittent or persistently gapped records. Missing observations need not be sensor failure: dry channels can change the estimand. Measurement definitions, quality flags and aggregation limit comparability. New histories are European mountain catchments, with few calibration/test components and no within-bundle donor variation. Regulation and groundwater status are unverified.

Calendar replay transfers missing shapes to complete years, not actual hidden observations or an identified MNAR mechanism. Derived and published events differ; incomplete discharge leaves some compound conditions unknown. A corrected snow label creates no new independent examples. Raw trajectory bands miss bias and transfer uncertainty; coverage is measured rather than assumed (SI Table S22).

Power projections stay low for five-point superiority when training/calibration variability is retained, even enlarging only the test set. We stop tuning rather than acquiring stations until a preferred representation wins. A future design should first improve on direct H30 uncertainty and complete historical diagnostics, then freeze predictive and calibration choices and collect independent hydrological and management-use evidence. A favorable subgroup or oracle cannot replace a failed primary comparison.

\section{Conclusions}

Incomplete temperature records should be assessed against an estimand and use, not one completeness percentage. Broad controls establish a substantial position effect while retaining length and reconstruction effects. The identity separates hidden-window error from compensating or harmful switching; oracle intervention does not demonstrate a large usable correction.

The proposed representation does not establish its predefined practical advantage over strong target-specific diagnostics. Unchanged transfer exceeds its nominal target at observed operating points; local recalibration and fitting are different claims. Screening can reduce loss relative to reconstructing all at some declared costs, but always-unknown answers can be cheaper. Unknown records, selected-sample distortion, current-vintage hindcast inputs and small independent cohorts remain part of the result. These distinctions define the present contribution and what a genuinely better design must demonstrate.

\section*{Open Research Statement}

Public products are reused under provider terms. Existing rights and measurement definitions are in \texttt{DATA\_RIGHTS.md}. New contracts, hashes, acquisition records, event/forcing resources and identity exclusions are under \texttt{results/wrr\_submission\_ready/major\_revision\_20261020c/}. The mountain resource has DOI \texttt{10.4211/hs.7a45ee06851e41d99e4affb7110d5f97}; continuous temperatures were reacquired, not inferred from extreme values. Local availability is not permission to redistribute every raw series.

Scripts 348--366 implement protocol, acquisition, recovery, calibration, transfer, references and power; scripts 367--371 and 373 generate final evidence, figures, tables, checks and corrected descriptive strata. Method and deployment locks identify fitted objects and thresholds before test scoring. The protocol and dated nonpredictive erratum are under \texttt{docs/reviews/wrr\_major\_revision\_20261020c/}. Script 368 generates figure sources; script 367 binds numerical macros to result rows and hashes. Previous manuscripts remain superseded history, not current evidence. A permanent archival deposit is required before submission.

\section*{Acknowledgments}
We thank the monitoring agencies and authors of the reused mountain-river resource. Funding, contributions, affiliations and corresponding-author details need verified author input, not inference from repository aliases.

